\scanpage{19}
\mspage{16}

Je dis que c'est un \uline{ordre}, i.e. qu'on ne peut avoir $x<y$ et $y<x$ (en effet, on aurait $x<x$, ce qui est faux, car $x<y$ implique $x\unsure{=|\ne}y$). Les segments sont les $[x,y]$, pour $x<y$, et pour \unsure{tt} \inserted{tel} segment, son origine est $x$ (pour l'orientation choisie \inserted{son extrémité $y$}). Ainsi, la structure de modèle \inserted{connexe} orienté de $1$-forme régulier est connue quand on connaît la structure d'ordre. Celle-ci \deleted{est} satisfait les conditions

elle est \inserted{str.} filtrante croissante (i.e. $\forall x,y\in L$, $\exists z\in L$ t.q. $x,y<z$) et str. filtrante décroissante, elle est \inserted{str.} loc. filtrante décroissante et loc. filtrante \deleted{\unsure{str.}} croissante : (si $\deleted{\unsure{x<}}\inserted{x\leq}y,z$, $\exists u$ tel que $x\deleted{\unsure{\leq}}\inserted{<}u\deleted{\unsure{\leq}}\inserted{<}y,z$, et forme duale)

\margin{donc pas de plus grand ni de plus petit \unsure{élém$^{\mathrm t}$}}

\margin{implique que c'est divisible, i.e. si $x<y$, $\exists z$ avec $x<z<y$}

Prouvons p.ex. qu'elle est \unsure{str.} filtrante croissante.

\mspage{18}

\rotatebox{90}{\reflectbox{\unsure{EXTRA STRONG}}}
\rotatebox{90}{\reflectbox{\unsure{EXTRA STRONG}}}
